\documentclass[11pt]{article}
\usepackage[margin=25mm]{geometry}
\usepackage{fontspec}
\setmainfont{Latin Modern Roman}
\usepackage{amsmath,amssymb,mathtools}
\usepackage{unicode-math}
\setmathfont{Latin Modern Math}
\usepackage{xurl,hyperref}
\hypersetup{colorlinks=true,urlcolor=blue}
\setcounter{secnumdepth}{0}
\setlength{\parindent}{0pt}
\setlength{\parskip}{5pt}
\setlength{\emergencystretch}{3em}
\newcommand{\sref}[2]{\hyperlink{source:#1:#2}{[S#2]}}
\newcommand{\cataloguescope}[1]{\paragraph{Recorded target.}#1}
\newcommand{\R}{\mathbb R}
\newcommand{\C}{\mathbb C}
\newcommand{\N}{\mathbb N}
\newcommand{\Z}{\mathbb Z}
\newcommand{\Q}{\mathbb Q}
\newcommand{\F}{\mathbb F}
\newcommand{\eps}{\varepsilon}
\begin{document}
\section{A non-polynomial complex potential with entirely real spectrum}\label{statement}

\subsection{Definitions and mathematical statement}

\paragraph{Shared notation.}$\mathbb N=\{1,2,\ldots\}$; the standard number fields have their usual meanings. All additional conventions are stated locally.


\paragraph{Definitions and notation.}A Schrödinger realization is specified by its potential, domain and boundary conditions. On $\mathbb R_+$ a Robin-type boundary condition can be written $y(0)\cos\theta+y'(0)\sin\theta=0$. A Riesz basis is the image of an orthonormal basis under a bounded invertible operator. An eigensystem may involve generalized eigenvectors when the operator is not diagonalizable. Augmented eigenfunctions in the source are spectral functionals used to represent evolution problems.

\paragraph{Catalogue wording retained for traceability.}
The record bundles source items 14.1--16.1.
\begin{enumerate}
\item Find a non-self-adjoint Schrödinger operator on $L^2(\mathbb R)$ or $L^2(\mathbb R_+)$ with a nonpolynomial potential, infinitely many real eigenvalues and at most finitely many non-real eigenvalues. The source compares with $H=-d^2/dx^2+x^m+a_1x^{m-1}+\cdots+a_m$ on $\mathbb R_+$ and the boundary condition $y(0)\cos\theta+y'(0)\sin\theta=0$, $a_j,\theta\in\mathbb C$.
\item Let $A_0=-d^2/dx^2+x^2$ on $L^2(\mathbb R)$ and $A=A_0+V$ for complex $V\in L^\infty(\mathbb R)$. Is its eigensystem a Riesz basis? The source records an affirmative result for $V\in L^p(\mathbb R)$, some $1\le p<\infty$.
\item Find additional applications for augmented eigenfunctions, and develop a spectral theory using them, beyond the two-point differential-operator initial-boundary representations described in the source.
\end{enumerate}
\paragraph{Scope and source qualification.}The inherited solved label is not fresh verification that all three source items have been settled. The title corresponds to only the first item, and no proof endorsement is supplied. The polynomial comparison and known basis case are historical source context. The catalogue title/status was paired with a machine-generated research attempt; that attempt is not peer-reviewed evidence of settlement. The statement covers the original source questions rather than only the title’s subcase.
\subsection{Short English statement}
Find nonpolynomial complex potentials with almost entirely real spectra, determine a harmonic-oscillator basis property, and develop augmented-eigenfunction spectral theory.
\subsection{Research attempt}
\paragraph{Model, date and investigated scope.}
OpenAI GPT-6 Astra (configured model \texttt{gpt-6-astra}), ultra reasoning; generated 11 October 2026. Constructed a nonpolynomial complex potential with explicitly proved infinitely many real eigenvalues, but did not exclude additional nonreal spectrum. Riesz-basis and augmented-eigenfunction items remain unproved; source-solved status is preserved.

\paragraph{Sources and literature check.}
Read the supplied complete definition, catalogue wording, scope qualifications and source list for this canonical ID. The original-source citations are inherited as supplied references; no fresh primary-source or novelty search was performed for this attempt. No current open-status claim is inferred from this calculation.

\paragraph{A Darboux construction with a deliberately incomplete spectral claim.}Fix real \(c\ne0\) and write \(z=x+ic\). Set \(u(z)=ze^{z^2/2}\); direct differentiation gives \((-d^2/dx^2+z^2)u=-3u\). Define \(A=d/dx-u'/u\), \(B=-d/dx-u'/u\). Multiplication of differential expressions yields
\[BA=H_0+3,\quad AB=H_1+3,\quad
H_0=-d^2/dx^2+z^2,\quad H_1=-d^2/dx^2+z^2-2+2/z^2.
\]
The potential of \(H_1\) is smooth on the real axis, nonpolynomial and genuinely complex. The algebra gives \(H_1A=AH_0\).
\paragraph{Infinitely many real eigenvalues.}Let \(\psi_0=e^{-z^2/2}\), and \(\psi_n=(-d/dx+z)^n\psi_0\). The commutator \([H_0,-d/dx+z]=2(-d/dx+z)\) and \(H_0\psi_0=\psi_0\) prove \(H_0\psi_n=(2n+1)\psi_n\). Each is a nonzero polynomial times the Gaussian, square integrable on the shifted horizontal line. Thus \(A\psi_n\) is also square integrable and satisfies \(H_1A\psi_n=(2n+1)A\psi_n\). It cannot vanish, since \(Af=0\) forces \(f\) proportional to the growing function \(u\). These functions lie in the natural maximal differential domain, since their derivatives and polynomial-potential products remain square integrable. Also \(1/u\) is square integrable and solves \(H_1(1/u)=-3/u\).
\paragraph{Remaining scope.}This explicitly produces infinitely many real eigenvalues, but the intertwining calculation alone has not excluded additional nonreal eigenvalues or proved a complete spectral correspondence of closed realizations. That is the first gap in the first subproblem. The bounded-perturbation oscillator Riesz-basis question and augmented-eigenfunction programme are entirely separate. The inherited source-solved label is retained without endorsing settlement of all three items; this reconstruction makes no new-solution claim.

\paragraph{Verification and adversarial audit.}
Checked factorization constants, absence of real-axis poles, ladder commutator, square integrability and nonzero transformed eigenfunctions; withheld completeness and all-bundle solution claims. This was an LLM adversarial audit, not independent human verification; no theorem prover was used. Prompts A then B applied to this ID: Checked factorization constants, absence of real-axis poles, ladder commutator, square integrability and nonzero transformed eigenfunctions; withheld completeness and all-bundle solution claims. Scope remains limited: Prove a complete spectral correspondence for the closed Darboux realizations to bound all nonreal eigenvalues.

Coordinator checked the Darboux factorization, shifted oscillator ladder and square-integrable transformed eigenfunctions, with spectral exhaustiveness expressly withheld. This additional check was model-assisted, not independent human verification.

\paragraph{Tools, process and contributions.}
Python 3 and local shell used to read supplied files and serialize this manuscript; mathematics derived symbolically in the text, with no numerical experiment or formal proof assistant. Contribution prompt A was followed by prompt B. GPT-6 Astra with ultra reasoning drafted the research and performed its stated self-audit; any additional coordinator contribution is identified in the verification and audit. A GPT-6-based Codex coordinator assembled the manuscript and checked repository integration; its exact variant and reasoning configuration were not exposed. No human mathematical contribution, independent human proof check or formal proof certification is supplied. Complete derivations and the consulted source references are included here; private reasoning transcripts are not published.

\paragraph{Final status and first remaining gap.}
\textbf{UNRESOLVED}. Prove a complete spectral correspondence for the closed Darboux realizations to bound all nonreal eigenvalues.

\section{Completion Estimate}
\noindent\textbf{COMPLETION: 0\%}
The calculation is an elementary benchmark or conditional reduction; it establishes no claimed new portion of the complete source target.

\subsection{Sources}
\begin{itemize}
\item[\hypertarget{source:statement:1}{S1}] ULAM.AI and the UnsolvedMath Contributors, UnsolvedMath, catalogue entry. Catalogue curation under CC BY 4.0; definitions and English exposition added locally.
\url{https://huggingface.co/datasets/ulamai/UnsolvedMath}.
\item[\hypertarget{source:statement:2}{S2}] American Institute of Mathematics, Mathematical aspects of physics with non-self-adjoint operators; workshop 8–12 June 2015, Problems 14.1–16.1.
\url{https://aimath.org/pastworkshops/nonselfadjointproblems.pdf}.
\end{itemize}
\label{end:statement}
\end{document}
